% Presentacion de tesis en Beamer. Estructura tomada de una plantilla de Overleaf;
% estilo propio en beamerthemeTesis.sty y colores en colores-tesis.sty.
% Compilar: latexmk presentacion.tex  (latexmkrc elige LuaLaTeX).
\documentclass[
  aspectratio=169,
]{beamer}

\usepackage[spanish, es-nodecimaldot, es-noshorthands, es-tabla]{babel}
\usepackage{csquotes}
\usepackage[backend=biber, style=ieee]{biblatex}
\addbibresource{../docs/referencias.bib}
\usepackage{booktabs}
\usepackage{tabularx}  % tablas con columnas X que se reparten el ancho
\usepackage{colortbl}  % \columncolor para sombrear una columna
\newcolumntype{Y}{>{\raggedright\arraybackslash}X}
% Celda de encabezado en tablas con fondo oscuro: texto blanco en negritas
% (\leavevmode evita que el color baje la celda una linea).
\newcommand{\encabezado}[1]{\leavevmode\color{white}\bfseries #1}
\usetheme{Tesis}
\usepackage{iconos-tesis}
\usetikzlibrary{calc}

% Marca temporal de cada lamina sin revisar. Se borra cuando ya no se use.
\newcommand{\pendiente}[1]{%
  \begin{tcolorbox}[colback=white, colframe=tenue, coltext=tenue, arc=4pt,
    boxrule=0.6pt, borderline={0.6pt}{0pt}{tenue, dashed}]
    \footnotesize\textbf{Pendiente de revisar.} #1
  \end{tcolorbox}}

\title[Detección de anomalías en redes de monitoreo]{Detección de anomalías en\\redes de monitoreo atmosférico}
\subtitle{Aprendizaje profundo con contraste espacio-temporal y exploración con imágenes sintéticas \textbf{TINTO}}
\author[J. E. Lili Beltrán]{Ing. José Emiliano Lili Beltrán}
\tutor{Dr. José Jaime Camacho Escoto}
\institute[PCIC · UNAM]{POSGRADO EN CIENCIA E INGENIERÍA DE LA COMPUTACIÓN · UNAM}
\date{Noviembre de 2026}
\subject{Tesis de maestría}
\keywords{detección de anomalías, redes de sensores, ozono, RAMA, LSTM}

\begin{document}

\begin{frame}[plain]
\titlepage
\end{frame}

% =====================================================================
\section{Introducción}

\begin{frame}{CONTEXTO DEL PROBLEMA}
\vfill
\centering
% Lamina 2 (contexto del problema): red RAMA, anomalias y amenazas, reto operativo.
% Ancho total: 13.8 cm (el ancho de texto de la lamina).
\begin{tikzpicture}[
    etiqueta/.style={font=\sffamily\footnotesize, text=texto, align=center, anchor=north, execute at begin node=\strut},
    titulo/.style={font=\sffamily\small, text=noche, anchor=west},
    flecha/.style={-{Latex[length=5pt]}, line width=1.2pt, ensX},
  ]
  % --- Titulos de los tres paneles ---
  \node[titulo] at (0.05,2.45) {{\color{acento}1.} Red RAMA (CDMX)};
  \node[titulo] at (3.95,2.45) {{\color{acento}2.} Anomalías y amenazas};
  \node[titulo] at (9.2,2.45) {{\color{acento}3.} Reto operativo};

  % --- 1. Red RAMA: estaciones y vecindad ---
  \coordinate (c) at (1.75,0.2);
  \fill[panel] (c) circle (1.72);
  \foreach \i/\a in {1/90, 2/162, 3/234, 4/306, 5/18}
    \coordinate (e\i) at ($(c)+(\a:1.15)$);
  \begin{scope}[ensR, dashed, line width=0.9pt]
    \foreach \i/\j in {1/2, 2/3, 3/4, 4/5, 5/1} \draw (e\i) -- (e\j);
    \foreach \i in {1,...,5} \draw (c) -- (e\i);
  \end{scope}
  \foreach \p in {c, e1, e2, e3, e4, e5} {
    \fill[white] (\p) circle (0.36);
    \pic[color=ensX, scale=1.05] at ($(\p)+(0,0.06)$) {antena};
  }
  % molecula de ozono
  \foreach \x/\y in {3.0/1.75, 3.25/1.95, 3.5/1.75}
    \fill[acento] (\x,\y) circle (0.13);
  \node[font=\scriptsize\bfseries, text=acento] at (3.25,1.42) {O\textsubscript{3}};

  % --- 2. Anomalias y amenazas: formato valido, valor falso ---
  \fill[panel, rounded corners=6pt] (3.95,-1.45) rectangle (8.85,1.75);
  \pic[color=ensX, scale=1.6] at (4.6,0.4) {antena};
  \node[draw=ensX, line width=1pt, fill=white, rounded corners=3pt,
        minimum width=1.3cm, minimum height=0.6cm,
        font=\footnotesize\bfseries, text=acento] (dato) at (6.3,0.3) {O\textsubscript{3} = 65};
  \draw[line width=1.2pt, ensX] (5.0,0.3) -- (dato.west);
  \draw[flecha] (dato.east) -- (7.65,0.3);
  \pic[color=ensX, scale=1.6] at (8.2,0.3) {servidor};
  \pic[color=acento, scale=1.3] at (6.3,1.25) {bicho};
  \draw[-{Latex[length=4pt]}, line width=1pt, acento] (6.3,0.9) -- (dato.north);
  \fill[noche] (dato.north east) circle (0.15);
  \pic[color=white, scale=1.1] at (dato.north east) {paloma};
  \node[etiqueta] at (4.6,-0.55) {sensor};
  \node[etiqueta] at (6.3,-0.55) {formato válido,\\valor falso};
  \node[etiqueta] at (8.2,-0.55) {servidor};

  % --- 3. Reto operativo: aviso al cierre del dia ---
  \fill[panel, rounded corners=6pt] (9.2,-1.45) rectangle (13.8,1.75);
  \pic[color=ensX, scale=1.6] at (10.0,0.3) {luna};
  \draw[flecha] (10.5,0.3) -- (11.05,0.3);
  \pic[color=ensX, scale=1.6] at (11.6,0.3) {servidor};
  \draw[flecha] (12.15,0.3) -- (12.7,0.3);
  \pic[color=ensX, scale=1.6] at (13.2,0.3) {llave};
  \fill[acento] (12.45,1.15) circle (0.33);
  \pic[color=white, scale=1.25] at (12.45,1.15) {campana};
  \node[etiqueta] at (10.0,-0.55) {cierre\\del día};
  \node[etiqueta] at (11.6,-0.55) {servidor};
  \node[etiqueta] at (13.2,-0.55) {revisar};
\end{tikzpicture}

\vfill
\note{%
\begin{itemize}
  \item En una red como la RAMA (32 estaciones en la CDMX; VERIFICAR: 31 evaluadas,
    33 en la matriz, y la red cubre también el Estado de México), el ozono varía mucho
    entre horas del día y entre zonas de la ciudad.
  \item Cuando hay una falla de sensor, una descalibración o un ataque de integridad en la
    transmisión, los datos llegan al servidor central dentro de rangos físicamente
    posibles, por lo que un umbral fijo no los detecta.
  \item El objetivo práctico es un aviso temprano en el servidor central, sin aumentar el
    costo de la red ni instalar nodos redundantes, que indique al cierre del día si una
    estación requiere revisión.
\end{itemize}}
\end{frame}

\begin{frame}
\begin{block}{HIPÓTESIS}
Es posible detectar, con un alto grado de confiabilidad, anomalías en la
información reportada en redes de sensores utilizando el historial reciente de
cada estación y el de sus estaciones vecinas sin incurrir en un exceso de
falsos positivos.
\end{block}
\begin{block}{OBJETIVO}
Diseñar un detector de anomalías basado en redes LSTM que combine la historia 
temporal de cada estación con la de sus estaciones vecinas, con alta capacidad 
de detección y una tasa de falsas alarmas acotada, en una red de sensores de ozono.
\end{block}
\end{frame}

\begin{frame}{TRABAJOS RELACIONADOS}
% Tabla comparativa: columnas = trabajos, filas = caracteristicas.
% Y = columna X de tabularx con texto alineado a la izquierda (definida en el preambulo).
% La ultima columna (este trabajo) va sombreada con \columncolor.
\centering
\footnotesize                          % letra mas chica solo en esta lamina
\renewcommand{\arraystretch}{1.12}     % aire entre filas
\setlength{\tabcolsep}{4pt}            % menos espacio entre columnas (normal: 6pt)
% Primera columna: 1.8 cm, en negritas y azul noche (\leavevmode evita que el
% color empuje la celda una linea hacia abajo).
\begin{tabularx}{\textwidth}{>{\leavevmode\bfseries\color{noche}}p{1.8cm} Y Y Y >{\columncolor{panel}}Y}
  \toprule
  % Encabezado: nombre, salto de linea dentro de la celda (\newline) y revista
  & PARAAD-SA
  & STC-RAAD
  & FedWLA
  & \textcolor{ensX}{\textbf{Este trabajo}}\\
  \midrule
  Modelo      & Autoencoder recurrente con atención; no supervisado
              & Autoencoder recurrente espacio-temporal; no supervisado
              & IDS federado; supervisado % VERIFICAR con el articulo
              & Ensamble de 5 LSTM; supervisado \\
  Usa vecinas & No & Sí & No aplica & Sí, con filtro de sospechosas \\
  Amenaza     & Datos: sesgo y deriva
              & Datos: anomalías de sensor % VERIFICAR
              & Red: tráfico malicioso
              & Datos: proporcional, escalón y sesgo \\
  Evaluación  & O\textsubscript{3}, 8 estaciones de Barcelona
              & Calidad del aire % VERIFICAR
              & X-IIoTID y Edge-IIoTset
              & O\textsubscript{3} horario, 31 estaciones RAMA \\
  \bottomrule
\end{tabularx}
\end{frame}

% =====================================================================
\section{Método}

\begin{frame}{MODELO DE AMENAZA Y EXPERIMENTOS}
% Datos tomados del protocolo (duracion: entero de 1 a 7 dias desde las 00 h) y de
% los registros previos de las revisiones 23, 23b y 24 (docs/registro_previo_rev2*.md).
% Encabezado con fondo azul noche (\rowcolor); filas blancas. Los \midrule de
% booktabs dejan una franja del color de fondo entre filas, como en la original.
\vfill
\centering
\small
\renewcommand{\arraystretch}{1.5}
\arrayrulecolor{fondo}  % las lineas entre filas, del color del fondo
\begin{tabularx}{\textwidth}{>{\leavevmode\raggedright\bfseries\color{noche}}p{4.2cm} Y >{\raggedright\arraybackslash}p{3.3cm}}
  \rowcolor{noche}
  \encabezado{Escenario} & \encabezado{Magnitudes evaluadas} & \encabezado{Duración} \\
  \rowcolor{white}
  1. Anomalía proporcional & $\pm\delta$, con $\delta$ = 20, 50 y 100\,\% & 1 a 7 días \\ \midrule
  \rowcolor{white}
  2. Escalón aditivo fijo  & $\pm$10, $\pm$20 y $\pm$40 ppb & 1 a 7 días \\ \midrule
  \rowcolor{white}
  3. Ataque coordinado a vecinas & 0, 1, 2 y 4 vecinas alteradas & 1 a 7 días \\ \midrule
  \rowcolor{white}
  4. Sesgo lento o sostenido & +20 y +40 ppb: escalón, rampa de 28 días y deriva de 56 días & 8 semanas (56 días) \\
\end{tabularx}
\vfill
\end{frame}

\begin{frame}{ARQUITECTURA DEL DETECTOR PROPUESTO}{Ensamble vecinal con filtro: ventana de 48\,h y decisión diaria}
\pendiente{Lámina 6: diagrama de la arquitectura (rehacer en TikZ o como imagen).}
\end{frame}

\begin{frame}{PROTOCOLO EXPERIMENTAL (2017–2026)}{Separación estricta de bloques anuales sin fuga de información}
\pendiente{Lámina 7: línea de tiempo; revisar la afirmación \enquote{sin fuga de información}.}
\end{frame}

% =====================================================================
\section{Resultados}

\begin{frame}{RESULTADOS · UNA ESTACIÓN ATACADA}{Sensibilidad (Recall) en anomalías proporcionales (\%), fijas (ppb) y despliegue}
\pendiente{Lámina 8: gráfica incompleta; las notas no coinciden con la gráfica.}
\end{frame}

\begin{frame}{ELECCIÓN DEL PUNTO DE OPERACIÓN}{Sensibilidad (Recall) frente a días con falsas alarmas al año}
\pendiente{Lámina 9: la gráfica no coincide con el título; decidir si va antes de los resultados.}
\end{frame}

\begin{frame}{RESULTADOS · ATAQUE COORDINADO A VECINAS (UMBRAL p99)}{Sensibilidad según el número de estaciones vecinas comprometidas (0 a 4)}
\pendiente{Lámina 10: el eje X salta de 2 a 4 con el mismo espacio.}
\end{frame}

\begin{frame}{RESULTADOS · SESGO SOSTENIDO DE LARGO PLAZO (UMBRAL p99)}{Horizonte de detección durante 8 semanas continuas}
\pendiente{Lámina 11: aclarar perfil de 14 días contra normalización de 30 días.}
\end{frame}

% =====================================================================
\section{Exploración con TINTO}

\begin{frame}{EXPLORACIÓN CON IMÁGENES SINTÉTICAS TINTO}{De una imagen diaria aislada a contexto temporal de 7 días}
\pendiente{Lámina 12: falta la tasa de falsas alarmas junto al recall.}
\end{frame}

% =====================================================================
\section{Referencias}

\begin{frame}[t, allowframebreaks]{REFERENCIAS}
\printbibliography[heading=none]
\end{frame}

\begin{frame}[plain]
\fondooscuro
\vfill
\centerline{\fuentetitular\LARGE\color{white}¡Gracias por su atención!}
\vfill\vfill
\end{frame}

% =====================================================================
% Respaldo: no cuenta en la numeracion "n / total".
\appendix
\section{Respaldo}

\begin{frame}{RESPALDO · TINTO}{Imagen TINTO de una estación-día}
\pendiente{Lámina 13: imagen sin título; decidir si se queda como respaldo.}
\end{frame}

\end{document}
